\subsection{Tipificación del incumplimiento de inversiones}

La separación de instrumentos requiere que el incumplimiento del programa de inversiones se discipline mediante una tipificación congruente con la obligación efectivamente incumplida. Una alternativa más directa consiste en sancionar la falta de ejecución oportuna y atribuible de una inversión previamente identificada, programada y financiada. Bajo este enfoque, la base económica correspondería a la parte no ejecutada por causas atribuibles a la empresa y el costo de la postergación se calcularía una sola vez para cada inversión, con independencia del número de metas o resultados a los que contribuya. Para ello, el programa de inversiones debe permitir identificar cada intervención, su monto programado, su cronograma y, cuando corresponda, su ámbito territorial.

Esta tipificación no excluye medidas correctivas o consecuencias específicas frente a incumplimientos autónomos de la calidad del servicio. La postergación de una inversión y la prestación por debajo de los estándares exigibles constituyen conductas diferentes, aunque puedan encontrarse relacionadas, por lo que cada una debe evaluarse mediante una consecuencia congruente con la obligación o afectación correspondiente. El Anexo~A desarrolla la identificación de la inversión no ejecutada, el cálculo de su costo postergado y el tratamiento de las inversiones vinculadas con varias metas.
